% ============================================================
% Capítulo: Marco Teórico (Bloque B.5)
% ============================================================
\chapter{Marco teórico}
\label{cap:marco-teorico}

Este capítulo sienta únicamente los fundamentos conceptuales
\textbf{estrictamente necesarios} para entender las decisiones y la
evidencia presentadas en el resto del documento -- no es un compendio
general de cada tema que toca este proyecto. Cuando un concepto ya se
desarrolla en profundidad en otro capítulo (la categorización completa
de los 46 requisitos, el detalle de las 13 ADR, el código real de cada
control OWASP), aquí se presenta solo el marco conceptual mínimo y se
remite explícitamente al capítulo que lo aplica.

\section{Ingeniería de requisitos: ISO/IEC/IEEE 29148:2018}
\label{sec:mt-requisitos}

La norma ISO/IEC/IEEE 29148:2018 define el vocabulario y el proceso de
referencia para la ingeniería de requisitos de sistemas y software,
incluyendo el patrón sintáctico ``[condición] [sujeto] \emph{shall}
[acción] [objeto] [restricción]'' que este proyecto usa como vara de
medida en su checklist de calidad
(\texttt{docs/checklists/incose2023-req.md}, aplicado en detalle en el
\autoref{cap:ingenieria-requisitos}). La elicitación de requisitos --
la actividad de descubrir, no solo registrar, lo que un sistema debe
hacer -- es un proceso empíricamente estudiado y propenso a errores
sistemáticos incluso en manos experimentadas: \citet{bano2019elicitation}
documentan, a partir de un análisis de errores comunes en entrevistas de
elicitación, que problemas como preguntas cerradas prematuras o la
proyección de soluciones antes de entender el problema son recurrentes
incluso en practicantes entrenados, y \citet{palomares2021elicitation}
confirman en un estudio de estado de la práctica en 12 empresas reales
que las técnicas de elicitación usadas en la industria rara vez siguen
un proceso formal completo, sino una combinación pragmática de entrevistas,
prototipado y retroalimentación iterativa -- exactamente el patrón que el
\autoref{cap:ingenieria-requisitos} documenta para SGB-SaaS. Este
capítulo no repite esa aplicación concreta; solo establece que la
elicitación de requisitos es, por naturaleza, un proceso falible que debe
verificarse después de ejecutado, no asumirse correcto por haberse
seguido un procedimiento nominal.

\section{Características de calidad de requisitos: INCOSE v4}
\label{sec:mt-incose}

El marco INCOSE v4 (Guide for Writing Requirements) define nueve
características que un requisito individual bien escrito debería
cumplir (Necessary, Appropriate, Unambiguous, Complete, Singular,
Feasible, Verifiable, Correct, Conforming) y seis características
adicionales a nivel de conjunto de requisitos (Complete, Consistent,
Feasible, Comprehensible, Able to be validated, Correct). Este capítulo
solo nombra el marco como fundamento conceptual; su aplicación completa,
requisito por requisito, con hallazgos y excepciones documentadas, está
en \texttt{docs/checklists/incose2023-req.md} y se resume en el
\autoref{cap:ingenieria-requisitos}.

\section{Modelo arquitectónico C4}
\label{sec:mt-c4}

El modelo C4 \citep{brown2023c4model} propone documentar la arquitectura
de un sistema de software en cuatro niveles de abstracción anidados --
Contexto, Contenedores, Componentes y Código -- cada uno pensado para una
audiencia y un nivel de detalle distintos, en vez de un único diagrama
monolítico que intenta servir a todas las audiencias a la vez. SGB-SaaS
usa los primeros tres niveles (el nivel de Código se considera cubierto
por el propio código fuente versionado). Este capítulo solo introduce el
modelo como notación; su aplicación concreta -- los tres niveles
completos, incluyendo el detalle de los 21 componentes del backend -- se
desarrolla en el \autoref{cap:diseno-arquitectura}.

\section{ISO/IEC 25010: modelo de calidad de producto de software}
\label{sec:mt-25010}

ISO/IEC 25010:2011 (SQuaRE) organiza la calidad de un producto de
software en ocho características de alto nivel -- Adecuación funcional,
Eficiencia de desempeño, Compatibilidad, Usabilidad, Fiabilidad,
Seguridad, Mantenibilidad y Portabilidad -- cada una con subcaracterísticas
propias, pensadas como un vocabulario común para discutir atributos de
calidad no funcionales de forma comparable entre proyectos. Este
capítulo solo introduce el marco conceptual; el mapeo completo de
SGB-SaaS contra las ocho características, con evidencia y prioridad
declarada para cada una, está en \texttt{docs/arquitectura/ISO25010.md}
y se transcribe en el \autoref{cap:diseno-arquitectura}.

\section{Principios REST}
\label{sec:mt-rest}

El estilo arquitectónico REST (\emph{Representational State Transfer})
describe un conjunto de restricciones para servicios web -- interfaz
uniforme, comunicación cliente-servidor sin estado (\emph{stateless}),
cacheabilidad explícita, y direccionamiento de recursos mediante URIs --
que SGB-SaaS adopta para el diseño de su API backend, tal como lo
implementa el framework usado en este proyecto \citep{walls2022springinaction}
mediante controladores anotados (\texttt{@RestController}) que exponen
recursos (\texttt{/api/v1/libros}, \texttt{/api/v1/prestamos}, etc.) sobre
los verbos HTTP estándar. La restricción de ausencia de estado en el
servidor es la que motiva, en particular, que la sesión del usuario
autenticado se transporte completa en cada petición (vía el token JWT,
\autoref{sec:mt-jwt}) en vez de mantenerse en memoria del servidor entre
peticiones.

\section{Autenticación con JWT y transporte seguro}
\label{sec:mt-jwt}

JSON Web Token (JWT, RFC~7519) es un formato compacto y auto-contenido
para representar afirmaciones (\emph{claims}) entre dos partes,
estructurado en tres segmentos separados por puntos -- \emph{header},
\emph{payload} y \emph{signature} -- codificados en Base64URL y firmados
(en el caso de SGB-SaaS, con un secreto simétrico HMAC), de modo que el
servidor puede verificar la integridad y autenticidad del token sin
necesidad de una consulta a base de datos ni de mantener estado de sesión
en memoria -- coherente con la restricción \emph{stateless} de REST
(\autoref{sec:mt-rest}). Este mecanismo no está exento de riesgos de
diseño: \citet{bucko2023jwtbehavior} muestran que un esquema JWT ingenuo
es vulnerable a robo y reuso de token, y proponen reforzarlo con
historial de comportamiento del usuario; \citet{shatnawi2024pythonjwt}
documentan, en un estudio empírico comparativo de librerías JWT en
Python, defectos de seguridad concretos (advertencias de \emph{linting}
de seguridad) presentes incluso en librerías JWT de uso común; y
\citet{ahmed2019jwtscheme} proponen regenerar el JWT en cada petición
del cliente con un componente de marca de tiempo verdaderamente aleatorio
para reducir la ventana de reuso de un token robado. Frente a este
panorama de riesgos documentados, la decisión de diseño de SGB-SaaS de
transportar el JWT en una cookie \texttt{HttpOnly} (en vez de en
almacenamiento accesible por JavaScript, como \texttt{localStorage}) --
detallada con código real en el \autoref{cap:implementacion} -- responde
directamente a la clase de riesgo que esta literatura describe: un
token en una cookie \texttt{HttpOnly} no puede leerse ni exfiltrarse
mediante una inyección de \emph{script} en el cliente (XSS), que es
precisamente el vector de robo de token más citado en la literatura de
seguridad de JWT. El transporte de esa cookie sobre una conexión cifrada
sigue, a su vez, las guías oficiales de configuración TLS de
\citet{mckay2019tls} (NIST SP~800-52 Rev.~2).

\section{Controles OWASP Top 10:2021}
\label{sec:mt-owasp}

El OWASP Top~10:2021 es una lista de referencia, mantenida por la
comunidad de seguridad web OWASP y actualizada periódicamente, de las
diez categorías de riesgo de seguridad más críticas y frecuentes en
aplicaciones web (control de acceso roto, fallos criptográficos,
inyección, mala configuración de seguridad, fallos de identificación y
autenticación, fallos de registro y monitoreo, entre otras). Este
capítulo solo nombra el marco como referencia; la auditoría real
ejecutada contra el stack Docker de SGB-SaaS, control por control, con
evidencia reproducible (\texttt{curl}) y corrección aplicada a cada
hallazgo, está documentada en \texttt{docs/mediciones/sec/} y se
instrumenta mediante \texttt{scripts/owasp-audit.sh}
(\autoref{cap:materiales-metodos}).

\section{Patrones de acceso a datos: ORM, procedimientos almacenados y \emph{cache-aside}}
\label{sec:mt-acceso-datos}

Un \emph{Object-Relational Mapper} (ORM) traduce automáticamente entre el
modelo de objetos de un lenguaje orientado a objetos y el modelo
relacional de una base de datos, a costa de un control más limitado sobre
la sentencia SQL exacta que finalmente se ejecuta. Spring Data JPA,
usado en el backend de SGB-SaaS \citep{walls2022springinaction}, genera
esas sentencias automáticamente para operaciones CRUD simples de una sola
tabla. Sin embargo, ese mismo automatismo puede degradar el rendimiento
en operaciones que involucran múltiples tablas relacionadas o lógica
transaccional compleja: \citet{colley2018ormperformance} documentan
empíricamente, comparando configuraciones de un framework ORM líder
contra SQL manualmente optimizado, patrones de rendimiento indeseables
(\emph{anti-patrones}) que el uso no cuidadoso de un ORM introduce. Un
procedimiento almacenado (\emph{stored procedure}, SP), en contraste,
ejecuta lógica SQL directamente dentro del motor de base de datos, con
control explícito sobre el plan de ejecución. SGB-SaaS adopta una
estrategia híbrida -- ORM para CRUD elemental, SP para operaciones
multi-tabla transaccionalmente sensibles -- cuyo mecanismo concreto de
invocación segura (parámetros nombrados bindeados, no concatenación de
cadenas) se documenta con código real en el \autoref{cap:implementacion}.

El patrón \emph{cache-aside} (también llamado \emph{lazy loading} de
caché) consulta primero una caché de acceso rápido (en este proyecto,
Redis) y solo recurre a la base de datos relacional cuando hay un fallo
de caché (\emph{cache miss}), momento en el que el resultado se escribe
de vuelta en la caché para peticiones futuras. \citet{kusuma2017rediscaching}
miden empíricamente el efecto de este patrón -- separando el tipo de
objeto cacheado (contenido estático vía acelerador HTTP, resultados de
consulta vía Redis) -- documentando reducciones reales de carga sobre la
base de datos relacional bajo tráfico concurrente. La instancia concreta
de este patrón en SGB-SaaS (anotación \texttt{@Cacheable} sobre el
catálogo de libros, con invalidación explícita vía \texttt{@CacheEvict})
y su evaluación empírica bajo carga (k6, comparación estadística
cache\_caliente vs.\ cache\_frío) se desarrollan en el
\autoref{cap:diseno-arquitectura} y en el \autoref{cap:materiales-metodos},
respectivamente.
